% Generated draft: requires mathematical and source review.
\documentclass[11pt]{article}
\usepackage[margin=25mm]{geometry}
\usepackage{fontspec}
\setmainfont{Latin Modern Roman}
\usepackage{amsmath,amssymb,mathtools,unicode-math}
\setmathfont{Latin Modern Math}
\usepackage{enumitem,needspace,xurl,hyperref}
\setlength{\parindent}{0pt}
\setlength{\parskip}{5pt}
\setlength{\emergencystretch}{3em}
\newcommand{\N}{\mathbb{N}}
\newcommand{\Z}{\mathbb{Z}}
\newcommand{\Q}{\mathbb{Q}}
\newcommand{\R}{\mathbb{R}}
\newcommand{\C}{\mathbb{C}}
\newcommand{\F}{\mathbb{F}}
\newcommand{\A}{\mathbb{A}}
\newcommand{\PP}{\mathbb{P}}
\newcommand{\E}{\mathbb{E}}
\newcommand{\eps}{\varepsilon}
\newcommand{\dd}{\,\mathrm{d}}
\begin{document}
\section*{UnsolvedMath ID: 30001465}
\textbf{Source problem code:} catalogue entry
\subsection{Definitions and mathematical statement}
A group \(G\) is torsion-free if \(g^m=1\), for \(g\in G\) and an integer \(m\geq 1\), implies \(g=1\). Let \(BG\) denote its classifying space, and let \(RG\) denote the group ring, consisting of finite formal sums \(\sum_{g\in G}r_g g\) with coefficients in \(R\), with multiplication extending the group multiplication and the multiplication in \(R\). For a spectrum \(\mathbf E\), write
\[
H_n(BG;\mathbf E)=\pi_n\bigl(\Sigma^\infty(BG_+)\wedge\mathbf E\bigr),
\]
where \(BG_+\) is \(BG\) with a disjoint basepoint. The source uses \(\mathbf K_R\) for an algebraic \(K\)-theory spectrum of \(R\), and \(\mathbf L_R^{\langle-\infty\rangle}\) for the indicated decorated algebraic \(L\)-theory spectrum.

\textbf{Source assertion (Farrell--Jones conjecture, as recorded).} For every torsion-free group \(G\), every regular ring \(R\), and every integer \(n\in\Z\), are both assembly maps
\[
\begin{aligned}
A^K_{G,R,n}\colon H_n(BG;\mathbf K_R)&\longrightarrow K_n(RG),\\
A^L_{G,R,n}\colon H_n(BG;\mathbf L_R^{\langle-\infty\rangle})&\longrightarrow L_n^{\langle-\infty\rangle}(RG)
\end{aligned}
\]
isomorphisms?

\textbf{Unresolved qualifications.} The supplied record does not define ``regular ring,'' specify whether rings are commutative or unital, or give the precise regularity and finiteness hypotheses intended. It also omits the involution and any orientation data needed to specify the \(L\)-theory objects and their assembly map. The all-integer grading requires a specified extension of algebraic \(K\)-theory to negative degrees, but the record does not state its convention. Consequently, the displayed question preserves the source formulation while leaving these coefficient and theory conventions unresolved. The record names this the Farrell--Jones conjecture but does not explain the hypotheses identifying this \(BG\)-based formulation with a formulation using the family of virtually cyclic subgroups. No such equivalence is asserted here.
\subsection{Short English statement}
The recorded question asks whether, for every torsion-free group and every regular coefficient ring, both displayed assembly maps recover the algebraic \(K\)- and decorated \(L\)-groups of the group ring in every integer degree. The meaning of regularity and the required theory conventions need clarification. The supplied curation labels the problem partially solved; that label is not an independent verification.
\subsection{Sources}
\begin{itemize}
\item Supplied local UnsolvedMath registry record, permanent ID \(30001465\), title ``Farrell--Jones Conjecture for Torsion-Free Groups,'' source problem code \texttt{catalogue entry}, in the collection ``Oberwolfach Reports Open Problems.'' The record gives the proposed year as 2010 and identifies no individual proposer.
\item \emph{Non-positive Curvature and Geometric Structures in Group Theory} (2010), Oberwolfach Reports, pp.\ 1165--1224, DOI \texttt{10.4171/owr/2010/20}. Supplied source report and literature-source link, labeled ``original OWR report'': \url{https://doi.org/10.4171/owr/2010/20}. The supplied background describes Oberwolfach Reports as documenting mathematical workshops and their open-problem discussions.
\item Supplied literature-source label ``Bartels--Lueck--Reich, hyperbolic groups'': \url{https://arxiv.org/abs/math/0701434}.
\item Literature assessment dated exactly \texttt{2026-08-21}: the supplied curation labels the assertion ``partially solved'' and reports results for major classes, including hyperbolic groups, while describing the universal torsion-free-group and regular-ring assertion as conjectural. This is supplied dated curation derived from the OWR record, not an independently verified literature assessment. The source itself requests independent verification of time-sensitive status and cited claims.
\end{itemize}
\end{document}
